\documentclass[11pt]{article}
\usepackage[margin=1.05in]{geometry}
\usepackage{amsmath,amssymb,amsthm}
\usepackage{xcolor}\usepackage{url}\def\UrlBreaks{\do\-\do\.\do\_\do\/}\sloppy\usepackage[colorlinks=true,linkcolor=blue,citecolor=blue,urlcolor=blue]{hyperref}
\newtheorem{theorem}{Theorem}[section]
\newtheorem{proposition}[theorem]{Proposition}
\newtheorem{lemma}[theorem]{Lemma}
\newtheorem{corollary}[theorem]{Corollary}
\theoremstyle{definition}
\newtheorem{definition}[theorem]{Definition}
\newtheorem{remark}[theorem]{Remark}
\newtheorem{question}[theorem]{Question}
\newcommand{\N}{\mathbb N}\newcommand{\Z}{\mathbb Z}\newcommand{\R}{\mathbb R}
\newcommand{\Pizt}{\Pi^0_2}\newcommand{\Sizt}{\Sigma^0_2}\newcommand{\Pith}{\Pi^0_3}\newcommand{\Sith}{\Sigma^0_3}
\newcommand{\Sio}{\Sigma^0_1}\newcommand{\Pio}{\Pi^0_1}
\newcommand{\cl}{\operatorname{cl}}
\title{The arithmetical complexity of approximation, residual and rigidity\\ properties of finitely generated groups}
\author{Research note, 2026-08-26\\ \small (Cairn graph \path{group-approximation}; node identifiers in \path{typewriter})}
\date{}
\begin{document}\maketitle
\begin{abstract}
We locate the recognition problems for the standard approximation properties of groups (sofic, hyperlinear, operator-MF, LEF, weakly and linear sofic), the residual properties (residual finiteness, residual $p$-finiteness), amenability and Kazhdan's property (T) in the arithmetical hierarchy, on three syntaxes: finite presentations, two-generator recursive presentations, and pairs of words in Higman's universal finitely presented group. The organizing statement is a second-level Rice theorem: every \emph{local approximation property} that contains the residually finite groups and fails for some group is $\Pizt$-hard on two-generator recursive presentations and undecidable with non-r.e.\ complement on finite presentations. Soficity is $\Pizt$-complete on recursive presentations, unconditionally. Hyperlinearity is either trivial or $\Pizt$-complete, so a non-hyperlinear group exists iff hyperlinearity of finitely presented groups is undecidable. Residual finiteness is $\Pith$-complete already for two-generator recursive presentations, one full level above every approximation property and above amenability, which is $\Pizt$-complete by Kesten's criterion. Property (T) is $\Sio$-complete on finite presentations, so its failure cannot be certified. Several of the finite-presentation cells are open and are stated precisely.
\end{abstract}

\section{Syntaxes and conventions}
$\mathrm{P}_{fp}$ denotes the set of codes of finite presentations whose group has property $\mathrm P$; $\mathrm{P}_{rec2}$ the set of indices of two-generator recursive presentations (a program enumerating relators) whose group has $\mathrm P$; $\mathrm P_V$ the set of pairs of words $(u_1,u_2)$ in a fixed finite generating set of Higman's universal finitely presented group $V$ such that $\langle u_1,u_2\rangle$ has $\mathrm P$; $\mathrm P_{enum}$ the countably generated enumerated presentations. $\mathrm{INF}=\{e: W_e \text{ infinite}\}$ and $\mathrm{FIN}$ are $\Pizt$- and $\Sizt$-complete; $\mathrm{TOT}=\{e:W_e=\N\}$ is $\Pizt$-complete; $\{e:\forall x\ W_{g(e,x)}\text{ finite}\}$ is $\Pith$-complete for suitable computable $g$; $\mathrm{COF}$ is $\Sith$-complete; $\mathrm{REC}$ is $\Sith$-complete.

\section{Local approximation properties and the marked-group space}
\begin{definition}[\path{local-approximation-properties-are-marked-closed}]
A \emph{finite table} is a finite set with a distinguished element and a partial binary operation, up to isomorphism; $\mathrm{tab}_G(F)=(F\cup F\!\cdot\! F,\ (x,y)\mapsto xy,\ 1)$ for finite $F\subseteq G$. A \emph{model predicate} is a predicate $M(T,\varepsilon)$ on tables and rationals $\varepsilon>0$, monotone in $\varepsilon$ and passing to sub-tables. The associated \emph{local approximation property} is $P_M(G):\iff \forall F\subseteq G$ finite, $\forall \varepsilon>0$: $M(\mathrm{tab}_G(F),\varepsilon)$.
\end{definition}
Sofic (permutation models, Hamming), hyperlinear (unitary models, normalized Hilbert--Schmidt), operator-MF (unitary models, operator norm, fixed separation), LEF (exact finite models), weakly sofic and linear sofic are of this form.
\begin{theorem}\label{thm:closed}
For every model predicate $M$: $P_M$ is hereditary and local; it is closed in the space of marked groups; and every finitely generated group $G=\langle X\mid N\rangle$ failing $P_M$ has a finitely presented cover $\langle X\mid N_0\rangle\twoheadrightarrow G$, $N_0\subseteq N$ finite, failing $P_M$.
\end{theorem}
\begin{proof}[Sketch]
A failing table $\mathrm{tab}_G(F)$ with tolerance $\varepsilon$ is described by finitely many word equalities and inequalities among words of length $\le r$. Any marked group agreeing with $G$ on words of length $\le 4r$ carries an isomorphic table, hence fails $P_M$; this is openness of the failure locus. The finitely many equalities are consequences of finitely many relators $N_0$; in $\langle X\mid N_0\rangle$ they hold by construction and the inequalities hold because they hold in the quotient $G$. (Sofic and hyperlinear instances are kernel-checked: \path{Covers/TableCover}, \path{Covers/HyperlinearTableCover}.)
\end{proof}

\section{The second-level Rice theorem}
\begin{theorem}[\path{second-level-rice-theorem-for-local-approximation-properties}]\label{thm:rice}
Let $P=P_M$ be a local approximation property such that (H1) every residually finite group has $P$ and (H2) some finitely generated group fails $P$. Then: (a) $P_{fp}$ is $\Sio$-hard and its complement is $\Pio$-hard, so neither is decidable and $\neg P_{fp}$ is not r.e.; (b) $P_{rec2}$ is $\Pizt$-hard and $\neg P_{rec2}$ is $\Sizt$-hard; (c) $P_V$ is $\Pizt$-hard; (d) if $P_{rec}\in\Pizt$ then $P_{rec2}$ and $P_V$ are $\Pizt$-complete and $P_{fp}\in\Pizt\setminus\Pio$, $\neg P_{fp}\in\Sizt\setminus\Sio$.
\end{theorem}
\begin{proof}[Sketch]
By Theorem~\ref{thm:closed} there is a finitely presented $E\notin P$. (a) $P$ is a Markov property (positive witness the trivial group, forbidden subgroup $E$, heredity), so Adian--Rabin reduces the word problem to $P_{fp}$. (b) The parametric FIN/INF ray switch (\path{arbitrary-forbidden-seed-hereditary-property-switch}, Lean \path{ParametricRecursiveSwitchPresentation}) gives countably generated recursive $S_e$ with $S_e\cong 1$ for $e\in\mathrm{INF}$ and $S_e\cong E$ for $e\in\mathrm{FIN}$; the two-generator bidirectional bridge $B_2(-)$ has residually finite value on $1$ and contains its argument, so $Q_e=B_2(S_e)$ has $P$ iff $e\in\mathrm{INF}$. (c) Embed $\ast_e Q_e$ in $V$ by Mikaelian's explicit Higman embedding. (d) Hardness plus membership.
\end{proof}
Kernel-checked half: \path{HereditaryPropertySwitchCompleteness.enumeratedCodeProperty\_pi02Hard} proves (b) on enumerated presentations for every hereditary $P$ with a finitely presented forbidden seed.

\section{Soficity: an exact classification}
\begin{theorem}[\path{sofic-recognition-has-a-pi2-upper-bound}]\label{thm:sofup}
For a recursive presentation $\langle X\mid r_1,r_2,\dots\rangle$ of $H$, with $d$ the normalized Hamming distance,
$H$ is sofic iff for all $(W,N,k)$: some $w\in W$ is trivial in $H$, or there is a finite $Y$ and $\sigma:X\to\mathrm{Sym}(Y)$ with $d(\sigma(r_j),\mathrm{id})\le 1/k$ for $j\le N$ and $d(\sigma(w),\mathrm{id})\ge 1-1/k$ for $w\in W$. Hence $\mathrm{SOFIC}_{rec}\in\Pizt$, with a primitive recursive certificate checker.
\end{theorem}
\begin{proof}[Sketch]
($\Rightarrow$) restrict a sofic model of the ball of radius $n$ to generators; right-invariance gives $d(\sigma(u),\phi(\bar u))\le(|u|-1)\varepsilon$. ($\Leftarrow$) choose words $u_g$ for the ball; every trivial word of length $\le 4n$ has area $\le A$ from $r_1,\dots,r_N$, so conjugation invariance and subadditivity give defects $\le A/k$; separation is read off the nontrivial words $u_g^{-1}u_h\in W$. (Finite-presentation instance kernel-checked: \path{SoficRecognitionPi02.pi02\_isSofic}; the strict placement $\Pizt\setminus\Pio$ is \path{sofic\_fp\_strict\_position}.)
\end{proof}
\begin{theorem}[\path{sofic-recognition-two-generator-recursive-is-pi2-complete}]
$\mathrm{SOFIC}_{rec2}$ and $\mathrm{SOFIC}_V$ are $\Pizt$-complete; $\mathrm{SOFIC}_{fp}\in\Pizt\setminus\Pio$ and $\mathrm{NONSOFIC}_{fp}\in\Sizt\setminus\Sio$.
\end{theorem}
\begin{proof} Theorem~\ref{thm:rice} with (H2) the kernel-checked finitely presented nonsofic group (\path{Endpoint/MainResults.exists\_finitelyPresented\_nonsofic\_group}) and Theorem~\ref{thm:sofup}. \end{proof}

\section{The hyperlinear dichotomy}
\begin{theorem}[\path{approximation-property-recognition-dichotomy}]
Exactly one of: (A) every group is hyperlinear, and $\mathrm{HYP}_{fp},\mathrm{HYP}_{rec2},\mathrm{HYP}_V$ are decidable (everything); (B) $\mathrm{HYP}_{rec2}$ and $\mathrm{HYP}_V$ are $\Pizt$-complete, $\mathrm{HYP}_{fp}$ is $\Sio$-hard and in $\Pizt\setminus\Pio$, $\mathrm{NONHYP}_{fp}$ is $\Pio$-hard and not r.e. In particular, a non-hyperlinear group exists iff a finitely presented one exists iff $\mathrm{HYP}_{fp}$ is undecidable iff $\mathrm{HYP}_{rec2}$ is $\Pizt$-complete.
\end{theorem}
The same holds for weak soficity and linear soficity over any computable field, with membership in $\Pizt$ supplied respectively by Glebsky's product-quotient characterization (normal-closure membership in a direct product of finite groups is one bounded template in every coordinate) and by the rank metric.
\begin{remark}
An undecidability proof for $\mathrm{HYP}_{fp}$ would produce a non-hyperlinear group without exhibiting one, exactly as $\mathrm{MIP}^*=\mathrm{RE}$ refutes Connes' embedding conjecture; it must deliver a computable family $e\mapsto G_e$ with exact relations on the positive side, the perfect-completeness obstacle of the Taller--Vidick LCS hardness. Whether $\mathrm{HYP}_{fp}$ is $\Pizt$-complete in case (B) reduces to hyperlinearity of the positive branch of the MF-safe finite-presentation compiler, i.e.\ to quasidiagonality of its regular trace (\path{mf-compiler-positive-branch-is-hyperlinear}).
\end{remark}

\section{Residual finiteness is one level higher}
Let $F=F(a,b)$, $a_j=b^jab^{-j}$, and for $J\subseteq\Z$ let $A_J=\langle a_j:j\in J\rangle$ and $H_J=\langle F,s\mid [s,A_J]=1\rangle$. For $J\subseteq\Z_{>0}$ let $G_J=\langle a,t\mid [a,t^jat^{-j}]=1\ (j\in J)\rangle$ and $J^s=J\cup(-J)\cup\{0\}$. Write $\cl(S)$ for the closure of $S\subseteq\Z$ in the profinite topology of $\Z$: $i\in\cl(S)$ iff $(i+n\Z)\cap S\ne\emptyset$ for all $n\ge1$.
\begin{theorem}[\path{centralizing-hnn-over-conjugate-basis-is-rf-iff-index-set-closed}, \path{shift-raag-family-is-rf-iff-symmetrized-index-set-closed}]\label{thm:crit}
The profinite closure of $A_J$ in $F$ is $A_{\cl(J)}$; $H_J$ is residually finite iff $\cl(J)=J$; $G_J=A(\Gamma_J)\rtimes\Z$ with $\Gamma_J$ the graph on $\Z$ with edges $\{i,i+j\}$, $j\in J$, and $G_J$ is residually finite iff $J^s$ is closed. Both families are LEF and sofic for every $J$.
\end{theorem}
\begin{proof}[Sketch]
If $i\in\cl(J)$ and $U\le F$ has finite index with $A_J\le U$, then $b^n\in U$ for some $n$ and $a_i=b^{-kn}a_jb^{kn}\in U$ for $j\in J$, $j=i+kn$. Conversely, for $g\notin A_{\cl(J)}$ involving a letter $a_i$, $i\notin\cl(J)$, pass to $\langle a\rangle*\langle b\mid b^n\rangle$ with $(i+n\Z)\cap J=\emptyset$ and $n$ exceeding all index differences in $g$: the normal closure of $a$ is free on $a_0,\dots,a_{n-1}$ (Kurosh), the image of $A_J$ is a proper free factor missing $a_{i\bmod n}$, and virtually free groups are subgroup separable. For $H_J$: closed $\Rightarrow$ residually finite by separating pinches and mapping to $Q*_{\psi(A_J)}(\psi(A_J)\times\Z)$, virtually free; not closed $\Rightarrow$ the pinch $[s,a_i]$, $i\in\cl(J)\setminus J$, dies in every finite quotient. For $G_J$: a finite-order image of $t$ folds $d\in\cl(J^s)\setminus J^s$ onto an edge, killing $[a_0,a_d]$; closed index sets are separated through $A(\Gamma_{J,m})\rtimes\Z/m$. LEF: truncate the shift graph to one period; sofic: right-angled Artin kernel, sofic-by-amenable (Elek--Szab\'o).
\end{proof}
\begin{theorem}[\path{index-set-profinitely-closed-is-pi3-complete}]
$\{e: \cl(W_e)=W_e\}$ is $\Pith$-complete. Hardness: $J_e=\{0\}\cup\bigcup_x\{2^x(2j+1): j<|W_{g(e,x)}|\}$ is closed iff every row is finite, because the sets $D_x=2^x+2^{x+1}\Z$ are pairwise disjoint clopen subsets of $\hat\Z$ covering $\Z\setminus\{0\}$, an infinite row $x$ has the closure point $-2^x$, and a finite row is closed.
\end{theorem}
\begin{theorem}[\path{residual-finiteness-two-generator-recursive-is-pi3-complete}]
$\mathrm{RF}_{rec2}$ is $\Pith$-complete (two generators is optimal), and so is $\mathrm{RF}_V$; $\mathrm{RF}_{fp}\in\Pizt\setminus\Pio$. The family $G_{J_e}$ with rows $\{2^x(4j+1):j<|W_{g(e,x)}|\}$ (closure point $-3\cdot2^x$, stable under symmetrization) witnesses hardness; it is LEF and sofic throughout, so the level does not drop under those promises. Residual $p$-finiteness is $\Pith$-complete for every prime $p$ (\path{residually-p-two-generator-recursive-is-pi3-complete}).
\end{theorem}
Upper bound: residual finiteness of a recursive presentation is $\forall w\,[w=1\ \vee\ \exists$ finite quotient separating $w]$ where the inner statement is $\Sizt$ because a finite quotient must kill every enumerated relator ($\Pio$); on finite presentations the relator check is finite and the level drops to $\Pizt$. Since LEF is $\Pizt$ on the same syntax, residual finiteness and LEF, which coincide for finitely presented groups, are separated by one quantifier alternation on recursive presentations, and the family $G_{J_e}$ realizes the gap.

\section{Amenability, free subgroups, finite presentability}
\begin{theorem}[\path{amenability-two-generator-recursive-is-pi2-complete}]
$\mathrm{AMENABLE}_{rec2}$ and $\mathrm{METABELIAN}_{rec2}$ are $\Pizt$-complete and $\mathrm{CONTAINS}\text{-}F_2{}_{rec2}$ is $\Sizt$-complete. Upper bound: Kesten's criterion, $G$ amenable iff $\forall k\,\exists n$: the number of trivial words of length $2n$ over $S=X\cup X^{-1}$ is at least $(1-1/k)^{2n}|S|^{2n}$, a count certified by derivations. Hardness: $G_J$ is amenable iff metabelian iff $G_J\cong\Z\wr\Z$ iff $J=\Z_{>0}$ (a right-angled Artin group is amenable iff its graph is complete), and contains $F_2$ otherwise; reduce from $\mathrm{TOT}$. $\mathrm{AMENABLE}_{fp}\in\Pizt\setminus\Pio$ by Adian--Rabin with trivial witness and $F_2$ forbidden.
\end{theorem}
\begin{proposition}[\path{finite-presentability-two-generator-recursive-is-sigma2-hard}]
$\mathrm{FP}_{rec2}$ is $\Sizt$-hard and in $\Sith$: $G_J$ is finitely presentable iff $J$ is finite (Bieri--Strebel forbids an ascending HNN structure over a finitely generated subgroup of the right-angled Artin kernel). Finite presentability is convex on index sets for any family that only adds relators to a fixed base, so such families cannot reach $\Sith$; on enumerated presentations $\mathrm{FP}$, finite generation, finiteness and property (T) are $\Sith$-complete. Solvability of the word problem on $rec2$ is $\Sith$-complete, since the word problem of $G_J$ is decidable iff $J$ is.
\end{proposition}

\section{Property (T)}
\begin{theorem}[\path{kazhdan-property-t-of-finite-presentations-is-sigma1-complete}]
$\mathrm T_{fp}$ and $\mathrm T_{rec2}$ are $\Sio$-complete; failure of (T) is not semidecidable, and no algorithm approximates the spectral gap of a finitely presented group to within a fixed additive error.
\end{theorem}
\begin{proof}[Sketch]
Semidecidability is Ozawa's criterion ($\Delta^2-\lambda\Delta$ a sum of hermitian squares in $\R[\Gamma]$), rationalized by Netzer--Thom. For hardness, the Rabin group $K(w)$ attached to a word $w$ of a group with unsolvable word problem is trivial when $w=1$ (kernel-checked, \path{RabinConstruction}) and otherwise an amalgamated free product over a rank-two free subgroup, proper in both factors; a nontrivial amalgam has no fixed point on its Bass--Serre tree, so fails Serre's FA, and (T) implies FA (Watatani). Property (T) is not hereditary, so this is not a Markov-property argument: the negative branch is a splitting. For two generators use $SL_3(\Z)$ versus $SL_3(\Z)\times\Z$, which is two-generated because $SL_3(\Z)$ is perfect.
\end{proof}

\section{Ledger and open cells}
\begin{center}\small
\begin{tabular}{lll}
property & $fp$ & $rec2$ / $V$\\\hline
sofic & $\Pizt\setminus\Pio$ (complete: open) & $\Pizt$-complete\\
operator-MF & $\Pizt$-complete (Cairn) & $\Pizt$-complete\\
LEF & $=\mathrm{RF}_{fp}$ & $\Pizt$-complete\\
hyperlinear / weakly / linear sofic & trivial or $\Sio$-hard in $\Pizt\setminus\Pio$ & trivial or $\Pizt$-complete\\
amenable & $\Pizt\setminus\Pio$ (complete: open) & $\Pizt$-complete\\
residually finite, residually $p$ & $\Pizt\setminus\Pio$ (complete: RF-Higman) & $\Pith$-complete\\
property (T) & $\Sio$-complete & $\Sio$-complete\\
finitely presentable & -- & $\Sizt$-hard, $\Sith$
\end{tabular}
\end{center}
Open cells (\path{open-cells-arithmetical-complexity-of-group-properties}): $\Pizt$-completeness of $\mathrm{SOFIC}_{fp}$ (needs a sofic-safe finite-presentation compiler: centralizing HNN extensions over non-amenable, non-profinitely-closed Mikhailova-type subgroups), of $\mathrm{HYP}_{fp}$ in case (B) (quasidiagonality of the MF compiler's regular trace), of $\mathrm{RF}_{fp}$ (the residually finite Higman problem) and of $\mathrm{AMENABLE}_{fp}$ (every Higman-type compiler introduces free subgroups); $\Sith$-completeness of $\mathrm{FP}_{rec2}$; and whether some group fails weak or linear soficity.


\section{Crossing to finite presentations: one gate for all cells}
The MF-safe finite-presentation compiler (\path{mf-safe-finite-presentation-compiler}) is uniform in its finitely presented seed $D$, and on the $\mathrm{INF}$ branch its output does not involve the seed at all: it is a fixed rope $R_e$ built from free groups by finite direct products, one amalgam over a Mikhailova-type subgroup, and one further HNN extension.
\begin{theorem}[\path{finite-presentation-rice-criterion-via-the-fixed-positive-rope}]\label{thm:fpgate}
Let $P$ be a subgroup-hereditary isomorphism-invariant property such that $P(R_e)$ for every $e\in\mathrm{INF}$, and $\neg P(D)$ for some finitely presented $D$. Then $P_{fp}$ is $\Pizt$-hard and $\neg P_{fp}$ is $\Sizt$-hard; with a $\Pizt$ upper bound, $P_{fp}$ is $\Pizt$-complete.
\end{theorem}
\begin{proposition}[\path{mf-compiler-positive-branch-is-torsion-free}]
Every positive-branch output $R_e$ is torsion-free: each building block is free, a finite direct product of free groups, a subgroup of such, an amalgam over a common subgroup, or an HNN extension of a torsion-free base, and torsion in an amalgam or an HNN extension is conjugate into a vertex group.
\end{proposition}
\begin{corollary}[\path{torsion-freeness-of-finite-presentations-is-pi2-complete}]
$\mathrm{TORSIONFREE}_{fp}$ is $\Pizt$-complete and $\mathrm{HASTORSION}_{fp}$ is $\Sizt$-complete: apply Theorem~\ref{thm:fpgate} with $D=\Z/2$ and the $\Pizt$ normal form of \path{torsion-freeness-recognition-is-pi2-complete}. This is the first finite-presentation cell besides MF to be closed at the second level.
\end{corollary}
Theorem~\ref{thm:fpgate} also explains the cells that remain open. The rope is MF and torsion-free, but it contains free subgroups and is not residually finite, so for amenability and residual finiteness this route is provably unavailable and their completeness cells need a different mechanism; for soficity and hyperlinearity the cells are \emph{exactly} the questions whether the one explicit group $R_e$ is sofic, respectively hyperlinear (\path{sofic-safe-finite-presentation-compiler}, \path{mf-compiler-positive-branch-is-hyperlinear}).



\section{Further rows, and where the hierarchy stops}
Torsion-freeness is $\Pizt$-complete on recursive presentations and, through the gate of the previous section, on finite presentations; perfectness and every bounded nilpotency class or derived length are $\Pizt$-complete (seeds $\Z$ and $F_2$); solvability of the word problem is $\Sith$-complete; and inside Higman's universal group the three levels one, two and three are realized by property (T), amenability and residual finiteness. Thus each of the first three arithmetical levels is realized by a natural, classically studied property; a natural level-four property is open, the candidate being ``residually a finite $p$-group for some prime $p$'', whose hardness needs index sets whose pro-$p$ closedness is controlled independently for each prime.

For \emph{isomorphism} the arithmetical hierarchy stops being the right measure. On finitely generated recursive presentations, isomorphism is $\Sith$ --- two finite tuples of words, each relator condition an r.e.\ derivation search --- and $\Pizt$-hard against a single fixed target (the lamplighter, identified inside the shift family by amenability). On countably generated enumerated presentations it is $\Pith$-hard, hence \emph{not} $\Sith$: with $A_e=\bigoplus_{x\in S_e}\Z/2$ for a uniformly $\Sizt$ index set $S_e$, one has $A_e\cong\bigoplus_\N\Z/2$ iff $S_e$ is infinite. In fact the enumerated isomorphism problem is $\Sigma^1_1$-complete already for computable torsion-free abelian groups (Downey--Montalb\'an, \emph{J.\ Algebra} \textbf{320} (2008) 2291--2300).


\section{The fourth level}
Fix a computable enumeration of pairs (prime, index) and put $\mathrm{row}(i,k)=\{p_i^{\,n!}+c_i:\ \mathrm{start}_i\le n<\mathrm{start}_i+k\}$ with $\mathrm{start}_i\ge i$ and factorial shifts $c_i=M_i!$, $M_i\ge i$; let $J_e=\{0,1\}\cup\bigcup_i\mathrm{row}(i,k_i(e))$.
\begin{lemma}[selective visibility]
Fix a prime $q$ and $m\ge1$, and let $\lambda$ be the exponent of $(\Z/q^m)^\times$. For every $p\ne q$ and every $n\ge\lambda$ we have $\lambda\mid n!$, hence $p^{\,n!}\equiv1\pmod{q^m}$ --- a threshold uniform in $p$; and $c_i\equiv0\pmod{q^m}$ once $M_i\ge q^m$. So all but finitely many elements of $J_e$ are $\equiv1\pmod{q^m}$.
\end{lemma}
Consequently any sequence of distinct elements converges to $1$ in the pro-$q$ topology, so $J_e$ is pro-$q$ closed whenever every row over $q$ is finite; and $p$-adically each infinite row over $p$ converges to its own shift $c_i\notin J_e$, so $J_e$ fails to be pro-$p$ closed exactly when some row over $p$ is infinite. (By contrast, a row that is an arithmetic progression of common difference $p$ is \emph{dense} in $\Z_q$ for every $q\ne p$, since $p$ is invertible mod $q^m$: the obvious first attempt cannot work.)
\begin{theorem}[\path{residually-p-for-some-prime-is-sigma4-complete}]
Being residually a finite $p$-group for \emph{some} prime $p$ is $\Sigma^0_4$-complete for two-generator recursive presentations.
\end{theorem}
The same family with a universal outer quantifier gives that being residually a finite $p$-group for \emph{every} prime is $\Pi^0_4$-complete, so both halves of the fourth level are realized. The mechanism stops there: the index family supplies exactly two quantifier resources, the finiteness of one row and the two indices ranging over the rows, and a natural fifth level would need another topological parameter of the same group rather than another index. On finite presentations the same properties collapse one level ($\Sigma^0_3$, $\Pi^0_3$) and their completeness is open, the rope route being unavailable because the rope is not residually finite. Hence each of the first four arithmetical levels is realized by a natural property of finitely generated groups: property (T) at level one; sofic, hyperlinear (side (B)), operator-MF, LEF, amenable, torsion-free, perfect, commutative, trivial at level two; residual finiteness and residual $p$-finiteness at level three; residual $p$-finiteness for some prime at level four.

\section{Separability is the exact obstruction}
Let $K$ be residually finite, $L\le K$ any subgroup, and $\Gamma=\langle K,v\mid [v,L]=1\rangle=K*_L(L\times\Z)$.
\begin{theorem}[\path{centralizing-hnn-is-residually-finite-iff-edge-is-separable}]
$\Gamma$ is residually finite iff $L$ is closed in the profinite topology of $K$.
\end{theorem}
\begin{proof}[Sketch]
If $L$ is closed, separate the finitely many pinch letters of a Britton-reduced $g\ne1$ from $L$ by a finite quotient $q:K\to Q$ and map $\Gamma\to Q*_{q(L)}(q(L)\times\Z)$, a finite graph of groups with finite and virtually cyclic vertex groups over a finite edge, hence virtually free and residually finite; the image of $g$ stays Britton-reduced. Conversely, for any $\chi:\Gamma\to P$ finite, $U=\{f\in K:\chi(f)\in\chi(L)\}$ has finite index and contains $L$, so every closure point $k$ of $L$ has $\chi(k)\in\chi(L)$ and therefore $\chi([v,k])=1$, while $[v,k]\ne 1$ for $k\notin L$.
\end{proof}
\begin{theorem}[\path{enumerated-subgroup-separability-is-pi3-complete}]
Separability of a recursively enumerated subgroup of $F_2$ is $\Pith$-complete; finitely generated subgroups are always separable (M.\ Hall), so the complexity appears exactly when infinitely many generators may appear over time.
\end{theorem}
This also settles what the sofic question for the compiler's rope is \emph{not}: by \path{centralizing-hnn-sofic-via-regular-edge-centralizer}, a stable letter in the centralizer of the edge exists exactly when the finite quotient separates the tested element from the edge image, so for a finitely generated edge the quotient-model route is equivalent to separability of that edge; Mikhailova subgroups are finitely generated with undecidable membership, hence non-separable, and the route dies. Any proof that the rope is sofic --- or even LEF --- must use finite models whose restriction to the base is not a homomorphism onto a finite quotient.


\section{What the remaining cells are}
Two of the open cells now have a precise shape.
\begin{proposition}[\path{rf-fp-completeness-reduces-to-subgroup-separability-hardness}]
For a finite tuple $w$ in $K=F\times F$, the group $\langle K,v\mid [v,w_1]=1,\dots,[v,w_n]=1\rangle$ is finitely presented and, by the criterion above, residually finite iff $\langle w\rangle$ is separable in $K$. Hence $\mathrm{RF}_{fp}$ is $\Pizt$-complete as soon as separability of finitely generated subgroups of $F\times F$ is $\Pizt$-hard --- a question about the best-understood non-subgroup-separable group, and one asking for hardness of a set already known to be $\Pizt$.
\end{proposition}
The sofic and hyperlinear finite-presentation cells are properties of a single explicit group, the rope $R$ of \path{the-mikhailova-rope-object-card}: finitely presented, MF, torsion-free, not residually finite, containing free subgroups of rank two, with soficity, LEF and hyperlinearity open, and with quotient models provably unable to settle soficity. Anyone attacking those cells should attack $R$ rather than rebuild a compiler.

\section{Machine-checked statements}
The following are kernel-checked Lean~4 declarations on the \path{main} branch of the development (module names under \path{GroupApproximation/Computability/} unless noted).
\begin{itemize}
\item \path{HereditaryPropertySwitchCompleteness.enumeratedCodeProperty_pi02Hard}: for every hereditary property that holds for the trivial group and fails for one finitely presented seed, the property of enumerated presentation codes is $\Pizt$-hard (via the computable parametric switch \path{ParametricRecursiveSwitchPresentation}).
\item \path{SoficRecognitionSecondLevel.soficCode_pi02Hard}, \path{nonsoficCode_sigma02Hard}: soficity of enumerated presentations is $\Pizt$-hard, nonsoficity $\Sizt$-hard.
\item \path{SoficMicrostateNormalForm.isSofic_iff_forall_answers}, \path{SoficRecognitionPi02.pi02_isSofic}, \path{sofic_fp_strict_position}: the permutation void-challenge normal form, $\Pizt$ membership on finite presentation codes, and $\mathrm{SOFIC}_{fp}\in\Pizt\setminus\Pio$, $\mathrm{NONSOFIC}_{fp}\in\Sizt\setminus\Sio$.
\item \path{SoficEnumeratedPi02.sofic_enum_pi02Complete}, \path{nonsofic_enum_sigma02Complete}: $\Pizt$-completeness of soficity and $\Sizt$-completeness of nonsoficity on enumerated presentation codes---the first machine-checked exact arithmetical classification of a group approximation property.
\item \path{LEFEnumeratedPi02.lef_enum_pi02Complete}, \path{nonlef_enum_sigma02Complete}: the same exact classification for local embeddability into finite groups on enumerated presentation codes (exact permutation models, no metric).
\item \path{ArithmeticalLedgerEndpoint.arithmeticalLedgerHolds}: the whole ledger as one closed proposition, with a \path{#audit_closed_axioms} line certifying that it rests on no hypothesis.
\item \path{TrivialEnumeratedPi02.trivial_enum_pi02Complete}, \path{AbelianEnumeratedPi02.comm_enum_pi02Complete}: the folklore rows (triviality, commutativity) as exact $\Pizt$-completeness on enumerated presentation codes.
\item \path{RFPresentationPi02.rf_fp_strict_position}: $\mathrm{RF}_{fp}\in\Pizt\setminus\Pio$, $\mathrm{NONRF}_{fp}\in\Sizt\setminus\Sio$; \path{RFEnumeratedHardness.rfCode_pi02Hard}, \path{AmenableEnumeratedHardness.amenableCode_pi02Hard}, \path{HyperlinearEnumeratedHardness.hyperlinearCode_pi02Hard_of_exists} (conditional), \path{IsoInvariantSwitchHardness.kazhdanCode_pi02Hard} (property (T), via an isomorphism-invariant variant of the switch), \path{ElementaryEnumeratedHardness} (finiteness, triviality): second-level lower bounds on enumerated presentation codes.
\item \path{ProfinitelyClosedIndexSet.pi03Complete_closedIndex}: $\{e:\cl(W_e)=W_e\}$ is $\Pith$-complete, with the dyadic rows; this is the computability core of the $\Pith$ residual-finiteness theorems.
\item \path{HyperlinearUndecidabilityRoute.not_computablePred_iff_exists_nonhyperlinear_code} (and, independently, a peer session's \path{exists_not_isHyperlinear_iff_codeProperty_not_computable}): hyperlinearity of finite presentation codes is undecidable iff some code presents a non-hyperlinear group, given undecidability of the word problem. \path{Covers/HyperlinearTableCover.exists_finitelyPresented_cover_of_not_isHyperlinear} (peer session): finitely presented non-hyperlinear covers.
\end{itemize}
The group-theoretic halves (bridge, Higman embedding, right-angled Artin kernels, Bass--Serre, Kesten, Ozawa) are paper-level in the research nodes cited above.
\end{document}
